
The theoretical and methodological foundation of Deltx rests upon an the integration of four closely connected domains of software engineering and artificial intelligence: \textbf{AI-generated text and code detection}, \textbf{hierarchical software quality models}, \textbf{predictive time-series modeling}, and \textbf{Explainable Artificial Intelligence (XAI)}.

\section{AI-Generated Text and Code Detection}

The increased usage of \textbf{Large Language Models} has necessitated the urgent development of algorithms capable of distinguishing synthetic text and source code from human authorship. Early detection approaches based on techniques such as perplexity analysis, token likelihood estimation, and watermarking have been found to lose effectiveness significantly when faced with domain changes or adversarial paraphrasing. \parencite{diversity_boosts_ai_generated_text_detection, why_ai_generated_text_detection_fails}.

Recent research highlights the need to move beyond surface-level structural patterns and focus more on deeper content-level linguistic, semantic, and surprisal-based features. The \textbf{DivEye framework} demonstrates that human-authored text exhibits significantly higher variability in lexical and structural unpredictability (surprisal) compared to the highly structured, mathematically optimized outputs of LLMs\parencite{diversity_boosts_ai_generated_text_detection}. Furthermore, foundational research on the \textbf{Global Word Distribution} illustrates that AI models converge tightly to global probability distributions resulting in distinct, quantifiable Letter Distribution Signatures\parencite{beyond_perplexity_character_distribution_signatures_mdta}. 

In the specific context of source code, advanced models like \textbf{EX-CODE} uses language model token probabilities to provide explainable classifications of synthetic code snippets, successfully uncovering the specific syntactic features that differentiate human-written from AI-generated logic\parencite{excode_robust_explainable_model_detect_ai_generated_code}.

\section{Hierarchical Software Quality Models}

Defining and measuring software quality in a consistent and practical manner has remained a long-standing challenge within Software Engineering. A major step toward standardization was introduced through the \textbf{ISO/IEC 25010 framework}, which describes software quality across multiple interconnected dimensions such as Functional Suitability, Performance Efficiency, Security, and Maintainability, each further divided into more specific sub-characteristics.\parencite{isoiec25010} However, bridging the gap between this theoretical framework and the operational metrics generated by static analysis tools like SonarQube is highly complex. 

Methodologies such as the \textbf{Squale (Software QUALity Enhancement)} model and the \textbf{Quamoco project} attempt to resolve this by introducing hierarchical aggregation techniques\parencite{quamoco_product_quality_modelling_assessment}. Squale, in particular, utilizes distinct mathematical formulas to normalize highly divergent metric values into a uniform scale, and then combine them using weighted, non-linear aggregation methods to produce higher-level quality indicators \parencite{computing_weighted_average_in_squale}. One of the key innovations of the Squale methodology is its ability to penalize severe outliers, ensuring that large number of minor or positive metrics cannot obscure the impact of a critical flaw. This helps overcome the aggregation limitations commonly found in simpler arithmetic-based evaluation models.\parencite{mord11a_continuous_weighted_metric_aggregation}.

\section{Predictive Time-Series Modeling}

In the domain of predictive modeling for software evolution, research has historically focused on predicting fault-proneness using classical statistical methods, such as \textbf{ARIMA}, or traditional machine learning algorithms like \textbf{Random Forests} and \textbf{Support Vector Machines}\parencite{software_service_supporting_software_quality_forecasting}. While \textbf{Long Short-Term Memory (LSTM)} networks were later introduced to capture temporal dependencies in code churn and metric evolution, more recent studies consistently demonstrate that \textbf{Transformer-based architectures} outperform LSTM networks in long-sequence, multivariate time-series forecasting tasks.\parencite{comparative_analysis_lstm_gru_transformer_construction_cost}. 

Specifically, the \textbf{PatchTST (Patch Time Series Transformer)} architecture has revolutionized the field. By segmenting continuous time series into discrete, subseries-level patches, PatchTST retains localized semantic temporal information while quadratically reducing the memory and computational complexity of the attention mechanism\parencite{decompatchtst_hybrid_framework}. Furthermore, its channel-independent design processes each multivariate feature as an individual univariate sequence that shares global weights. This approach has proven highly effective in reducing cross-metric noise within complex and heterogeneous datasets.\parencite{patchtst_nixtla}.

\section{Explainable Artificial Intelligence}

Finally, integrating deep learning architectures into risk assessment systems requires a high degree of transparency in order to build developer trust and support effective, actionable remediation. \textbf{Explainable AI (XAI)} frameworks, particularly \textbf{SHAP (SHapley Additive exPlanations)} and \textbf{LIME (Local Interpretable Model-agnostic Explanations)}, have become essential components of modern predictive systems\parencite{comparing_explainable_ai_models_shap_lime_electric_field}. 

\textbf{SHAP} utilizes cooperative game theory to calculate the precise marginal contribution of each input feature to the model's final prediction, providing both global and local interpretability with strict mathematical consistency\parencite{explainability_testing_ai_shap_lime_testriq}. Applying SHAP to multivariate software defect prediction allows engineers to trace forecasted quality degradations back to their exact historical drivers.

\vspace{0.5cm}

The synthesis of these four domains, AI detection, quality modeling, predictive forecasting, and explainability—provides the comprehensive theoretical foundation upon which Deltx is architected.
